% Full literature review.
% 42 papers across 5 sections; ~3800 words.
% Research question: How can deep learning methods for BSDEs with jumps be
% applied to solve Nash equilibria and detect tacit collusion in multi-agent
% dealer market models, specifically the Cont--Xiong (2024) intensity-controlled
% market-making game?

\section{Introduction}\label{sec:lit-intro}

Multi-agent dealer market-making sits at an uncomfortable
intersection in the mathematical literature.  The $N$-player
stochastic-differential-game formulation of \citet{cont2024}
admits, in principle, a probabilistic treatment via backward
stochastic differential equations (BSDEs) and mean-field games
(MFGs), and a numerical treatment via deep-learning methods.  In practice, no
existing framework simultaneously handles all three ingredients:
\citet{han2022} give a curse-of-dimensionality-free deep-BSDE
method for McKean--Vlasov FBSDEs but without jumps;
\citet{wang2023} extend deep BSDE to jump-diffusion processes but
without mean-field coupling; \citet{cont2024} solve their own
game empirically via multi-agent reinforcement learning but
without the equilibrium grounding or interpretable gradients that
a principled method would supply.  Combining all three ---
jump-augmented BSDEs, mean-field coupling, and deep-learning
numerics --- into a single computational method for a realistic
market-microstructure model remains open, and the failure modes
that emerge when one tries are not systematically documented in
the literature.

This review is framed around a single target problem: the multi-agent
dealer market-making game of \citet{cont2024}.  Every piece of theory
surveyed in \S\ref{sec:lit-math}, every numerical tool catalogued in
\S\ref{sec:lit-num}, and every empirical finding discussed in
\S\ref{sec:lit-empirical} is selected because it is needed to solve
that game.  \S\ref{sec:lit-math} assembles the probabilistic
scaffolding --- BSDEs, jump processes, mean-field limits, master
equations --- that lets us write down the Cont--Xiong value function.
\S\ref{sec:lit-num} covers the deep-learning numerical machinery that
makes the resulting high-dimensional mean-field BSDE with jumps
computable.  \S\ref{sec:lit-empirical} positions the thesis against
the multi-agent reinforcement-learning literature in which
Cont--Xiong was first solved empirically, and isolates what a
principled BSDE treatment adds.  \S\ref{sec:lit-conclusion}
identifies what no existing work yet provides.

The starting point is \citet{pardoux1990}, who established the
existence and uniqueness of adapted solutions to BSDEs under Lipschitz
conditions by a probabilistic fixed-point argument.  Their result
provides the theoretical bedrock on which every subsequent development
rests: it guarantees that a BSDE, despite being specified ``backwards''
from a terminal condition, has a well-defined solution adapted to the
forward filtration of a Brownian motion.  The translation of this
technology into quantitative finance was effected by
\citet{elkaroui1997}, who showed that pricing, hedging, and risk
management under constraints admit a natural BSDE formulation.
For a European contingent claim $\xi = g(S_T)$ in a
Black--Scholes market with risk-free rate $r$, drift $\mu$ and
volatility $\sigma$, the no-arbitrage price and replicating
strategy $(Y_t, Z_t)$ satisfy the linear BSDE
\[
  -dY_t \;=\; \bigl[r Y_t + Z_t \theta\bigr] dt - Z_t\, dW_t,
  \qquad Y_T \;=\; g(S_T),
\]
where $\theta = (\mu - r)/\sigma$ is the market price of risk and
$Z_t$ is the hedging portfolio in the risky asset.  Their
framework unified arbitrage-free pricing, superhedging via
\emph{reflected} BSDEs (for American options), and Epstein--Zin-type
recursive utilities into a single probabilistic language, turning
BSDEs from a purely theoretical object into a practical tool for
derivative pricing and for the study of incomplete markets.

In parallel with the BSDE literature, Lasry and Lions --- and
independently Huang, Caines, and Malham\'{e} --- introduced
mean-field games as a tractable approximation of $N$-player
stochastic differential games in the regime $N \to \infty$.
\citet{carmona2013} recast this analytic theory in purely
probabilistic terms, showing that an MFG equilibrium is
characterised by a McKean--Vlasov forward-backward SDE system in
which the distribution $\mu_t$ of a representative agent's state
enters both the drift and the terminal condition; the stochastic
maximum principle then supplies the natural optimality condition,
avoiding the infinite-dimensional PDEs of the analytic formulation.
Extensions have since covered games of optimal timing
\citep{carmona2017} and the long-time ergodic behaviour of MFG
systems \citep{cardaliaguet2012}; for a comprehensive survey of
analytical and numerical methods, see \citet{achdou2020}.

% ============================================================
\section{Mathematical Framework}\label{sec:lit-math}

Any probabilistic treatment of BSDEs rests on Brownian motion, the
It\^{o} integral, and the classical theory of stochastic differential
equations, standard references for which are \citet{oksendal1998}.
Within this setting, a BSDE is a pair of adapted processes $(Y_t, Z_t)$
satisfying
\[
Y_t \;=\; \xi + \int_t^T f(s, Y_s, Z_s)\,ds - \int_t^T Z_s\,dW_s,
\qquad t \in [0, T],
\]
where $\xi$ is an $\mathcal{F}_T$-measurable terminal condition and
$f$ is the driver.  \citet{pardoux1990} proved that whenever $f$ is
uniformly Lipschitz in $(y, z)$ with $\xi$ and $f(\cdot, 0, 0)$ in
$\mathbb{H}^2$ (the Hilbert space of square-integrable progressively
measurable processes), there exists a unique adapted pair
$(Y, Z) \in \mathbb{H}^2 \times \mathbb{H}^2$ solving the BSDE,
obtained as the fixed point of the contraction
\[
  \Phi(Y, Z)_t \;:=\; \mathbb{E}\!\left[\xi + \int_t^T f(s, Y_s, Z_s)\,ds \,\Big|\,
    \mathcal{F}_t\right]
\]
on the Banach space $\mathbb{H}^2 \times \mathbb{H}^2$.  A large
body of later work has explored the structural consequences of this
theorem.  The Pardoux--Peng existence theory has since been extended in
several directions relevant to this project.  \citet{feng2018}
established the weak-solution correspondence between quasi-linear
PDEs and FBSDEs, which applies to the Cont--Xiong HJB system whose
coefficients are quasi-linear rather than classical.  \citet{peng2004} showed
that BSDEs naturally generate \emph{nonlinear expectations}.  Given
a driver $g(t, z)$ independent of $y$, the $g$-expectation of a
terminal value $\xi \in L^2(\mathcal{F}_T)$ is defined as
\[
  \mathcal{E}^g[\xi] \;:=\; Y_0,
  \qquad \text{where } (Y, Z) \text{ solves the BSDE with driver }
  g \text{ and terminal } \xi.
\]
When $g$ is positively homogeneous and subadditive in $z$, the
functional $\xi \mapsto \mathcal{E}^g[\xi]$ satisfies the
Artzner--Delbaen--Eber--Heath axioms (monotonicity, translation
invariance, positive homogeneity, subadditivity) of a coherent
risk measure; convex $g$ yields a convex risk measure in the
sense of F\"ollmer--Schied.  This elevates BSDE theory to an
axiomatic foundation for risk measurement, placing the $Z$-process
learned by later deep-BSDE algorithms in the role of a
risk-sensitivity functional.  The crucial martingale representation result underlying
existence — that any square-integrable martingale adapted to a
Brownian filtration admits a representation as a stochastic integral —
is what makes the $Z$-process well-defined and is also what Han, E,
and Jentzen will later exploit to parameterise $Z$ by a neural
network.

Mean-field BSDEs (MF-BSDEs) extend this picture by allowing the
driver or the terminal condition to depend on the law of the
solution itself.  \citet{buckdahn2009viscosity} established the
connection between MF-BSDEs and non-local partial differential
equations via viscosity solution theory: the solution of an MF-BSDE
provides the unique viscosity solution of an associated non-local
PDE on the state space.  The non-locality reflects the fact that the
distribution $\mu_t$ of the forward process enters the PDE, which
immediately precludes any classical finite-difference numerical
scheme above a few dimensions.  A complementary paper
\citep{buckdahn2009limit} establishes the mean-field limit itself:
starting from an $N$-player FBSDE system with empirical measure
$\mu_t^N = N^{-1} \sum_{i=1}^N \delta_{X_t^{i,N}}$, the authors
prove that each $i$-th particle converges in $L^2$ to a
McKean--Vlasov BSDE with law $\mu_t$ at a rate
\[
  \sqrt{\mathbb{E} \sup_{t \leq T} \bigl|X_t^{i,N} - X_t^i\bigr|^2
  + \mathbb{E} \sup_{t \leq T} \bigl|Y_t^{i,N} - Y_t^i\bigr|^2 }
  \;\leq\; \frac{C}{\sqrt{N}},
\]
for a constant $C$ depending on the Lipschitz constants of the
drift, diffusion, and driver.  This rate is of direct relevance
to every empirical study of mean-field approximation at finite
$N$: it gives a precise, non-asymptotic benchmark against which
finite-$N$ simulation error can be compared.

The jump-augmented case has a parallel classical theory that
predates the deep-learning numerics discussed in \S\ref{sec:lit-num}.
\citet{tang1994} established necessary conditions for optimal
control of stochastic systems with random jumps, extending
Pontryagin's maximum principle to a setting in which execution events
are discrete.  \citet{barles1997} gave the BSDE/integro-partial
differential equation correspondence that is the direct analogue of
the Pardoux--Peng--Feynman--Kac bridge for Brownian BSDEs: the
solution of a BSDE with a compensated Poisson jump term is the
unique viscosity solution of an associated IPDE.  Together these
papers provide the probabilistic foundation on which any numerical
jump-BSDE scheme --- including the deep BSDEJ method of \S
\ref{sec:lit-num}, and any MV-fictitious-play BSDEJ solver built
on top of it --- must ultimately rest.

The definitive probabilistic treatment of MFGs is the two-volume
monograph of \citet{carmona2018}.  Volume I develops the mean-field
FBSDE framework, the stochastic maximum principle, and
well-posedness under structural conditions.  Volume II extends the
theory to mean-field games with common noise — where all agents are
exposed to a shared stochastic shock, inducing conditional
distributions $\mu_t(\omega^0)$ rather than deterministic measure
flows — and to the master equation on the Wasserstein space.  The
master equation is a PDE in $(t, x, \mu)$ whose solution encodes the
entire MFG value function, and its probabilistic study combines the
Lions derivative with the FBSDE-flow machinery.  Existence in the
common-noise case was established independently by
\citet{carmona2016} using strong and weak formulations and a
Schauder fixed-point argument.  Classical solutions of the master
equation, and hence the differentiability of the MV-FBSDE flow, were
obtained by \citet{chassagneux2014} under strong
regularity assumptions on the coefficients.

On the control-theoretic side, \citet{bayraktar2018}
obtained a randomised dynamic programming principle and a nonlinear
Feynman--Kac representation for McKean--Vlasov value functions,
realising them as the $Y$-component of FBSDEs with autonomous
forward processes.  This result is the direct analogue, in the
mean-field setting, of the classical Feynman--Kac bridge that
motivates the deep BSDE method.  \citet{wuzhang2020} sharpened
the master-equation theory by introducing a genuine viscosity-solution
notion for parabolic master equations on the Wasserstein space, with
closed-loop controls, thereby removing some of the regularity
restrictions that Chassagneux et al.\ needed.

Finer properties of the mean-field limit were studied by
\citet{delarue2019}, who established a central limit theorem for
the $N$-player game fluctuations around their mean-field limit and
characterised the Gaussian correction in terms of a linearised
master equation.  Their CLT result is particularly relevant for the
present thesis because the Cont--Xiong market-making game is
naturally formulated in a finite-$N$ regime, and the $O(1/\sqrt{N})$
propagation of chaos of \citet{buckdahn2009limit} combined with the
CLT corrections of Delarue et al.\ gives a quantitative picture of
how finite-$N$ Nash equilibria deviate from mean-field equilibria.
Finally, \citet{zhang2023} studied stochastic recursive
control problems of McKean--Vlasov type via the dynamic programming
principle on an infinite-dimensional HJB equation, providing a
further link between MV-FBSDE theory and the recursive-utility
literature.  Together, these results yield a mathematical framework
that is well-posed, quantitatively understood at finite $N$, and
equipped with a viscosity and a probabilistic representation — but
still without a computable solution in any high-dimensional setting.

% ============================================================
\section{Numerical Methodology}\label{sec:lit-num}

The numerical analysis of BSDEs follows two historical phases.
The first, beginning with \citet{kloeden1992, kloeden1994} for
forward SDEs, establishes the discretisation benchmarks: for a
step size $\Delta t$ the Euler--Maruyama scheme has strong
convergence order $1/2$ (i.e.\
$\mathbb{E}|X_T - X_T^{\Delta t}| = O(\sqrt{\Delta t})$) and weak
order $1$, while the Milstein correction attains strong order $1$
at the cost of requiring derivatives of the diffusion
coefficient.  \citet{pengxu2011} extended these ideas to BSDEs,
developing backward time-stepping schemes --- implicit Euler with
rate $|Y_0^{\Delta t} - Y_0| = O(\sqrt{\Delta t})$ under standard
smoothness and Crank--Nicolson with rate $O(\Delta t)$ --- coupled
with conditional-expectation estimators (quantisation, Malliavin
weights, least-squares regression).  These methods are rigorous
and enjoy classical convergence rates, but they all share a
fatal limitation: computing the conditional expectations
$\mathbb{E}[Y_{t_{n+1}} | \mathcal{F}_{t_n}]$ by grid-based or
regression methods scales as $O(M^d)$ in the state dimension $d$,
so that $d \gtrsim 4{-}5$ is effectively out of reach.  For the
dealer-market problem of \citet{cont2024} this renders the
direct approach infeasible already at moderate numbers of
competitors.

The second phase is inaugurated by \citet{ehanjentzen2017} and its
journal version \citet{han2018}, who reformulate
a parabolic PDE as a BSDE and parameterise the gradient process $Z$ by
a deep neural network, with $Y_0$ being a single learnable scalar.
The training objective is the mean-square deviation of the simulated
terminal value $Y_T^\pi$ from the prescribed $g(X_T)$.  Once trained,
the network encodes both the value function $u(0, x) = Y_0$ and its
gradient $\nabla u(t, x) = Z_t/\sigma$ along all simulated paths.
Empirically, this approach handles up to 100 dimensions on a laptop;
the authors reported successful experiments on the Black--Scholes
PDE with correlated assets, the Hamilton--Jacobi--Bellman equation
for a linear-quadratic controller, and the Allen--Cahn PDE.  The
extension to fully nonlinear second-order BSDEs was carried out by
\citet{beck2019}, who treated 100-dimensional
Black--Scholes--Barenblatt, HJB, and $G$-Brownian-motion problems
under the 2BSDE framework, with neural networks approximating both
the gradient and the Hessian.  Both of these papers are empirical:
rigorous convergence proofs for the original deep BSDE method are
still absent, a gap that \citet{carmona2019} partially
fill by giving two algorithms for ergodic McKean--Vlasov control
and proving convergence under network-architecture-dependent rates.

The step from BSDEs to McKean--Vlasov FBSDEs introduces a second
source of difficulty: the coefficients now depend on the unknown law
$\mu_t$ of the forward process, so the training loop must couple
neural-network learning with a distribution-matching fixed-point
iteration.  \citet{guo2023} develop a general mean-field game (GMFG)
framework in which the coupling passes through the joint
state--action distribution $(\mu_t, \nu_t)$ rather than the state
distribution alone: rewards $r(x, a, \mu, \nu)$ and transition
kernels $P(\cdot | x, a, \mu, \nu)$ depend on the empirical action
law as well as the empirical state law, generalising classical
MFGs to McKean--Vlasov type.  They prove convergence of
fictitious play to the GMFG Nash equilibrium under Lipschitz and
contraction conditions on the associated Bellman operator,
handling $\mu_t$ computationally through empirical particle
representations.  The decisive theoretical advance for the
present thesis is \citet{han2022}, who propose a novel
deep-learning algorithm for MV-FBSDEs with general mean-field
interactions and prove that its convergence is \emph{free of the
curse of dimensionality} under monotonicity conditions.  Their proof
uses a class of integral probability metrics and a fictitious-play
decomposition: the outer loop updates the mean-field measure using
a population-averaged surrogate, and the inner loop solves a
standard FBSDE by deep BSDE techniques.  The complexity is shown to
grow only polynomially with dimension, so long as the Lasry--Lions
monotonicity condition holds on the coupling term.  This
monotonicity requirement is a real constraint in applications: for
the Cont--Xiong coupling, global Lasry--Lions monotonicity on the
Wasserstein space has not been established, and empirically
monotonicity can be verified only along a one-dimensional slice
parameterised by the mean competitor quote.  The Han--Hu--Long
convergence guarantee therefore applies only on the restricted
slice of measure space relevant to a coupling like Cont--Xiong's,
a subtlety that the original paper does not address.  Modulo that
caveat, the Han--Hu--Long algorithm is the natural template on
which any MV-fictitious-play BSDEJ solver for the Cont--Xiong
game would be built.

The extension of deep BSDE to jump-diffusion processes and forward
BSDEs with jumps (FBSDEJs) is due to \citet{wang2023}, whose work
is the computational counterpart of the classical BSDE/IPDE
equivalence of \citet{barles1997}.  Wang and collaborators
parameterise the jump coefficient $U_t^a, U_t^b$ alongside $Z_t$ by
separate neural networks and report convergence in up to 100
dimensions.  Two caveats warrant emphasis.  First, as of writing,
\citet{wang2023} has accumulated only a handful of citations,
making the deep BSDEJ numerical literature considerably less
mature than the Brownian-only deep BSDE literature it extends;
any downstream application must therefore do its own due
diligence on convergence, on architectural ablations, and on the
exact form of the compensated-martingale correction, which is
subtle to implement correctly.  Second, the Cont--Xiong model as currently formulated
has discrete inventory $q_t \in \{-Q, \ldots, Q\}$, so the
associated BSDE is pure-jump: the diffusion coefficient $Z_t$
vanishes identically, and the network learns only
$(U_t^a, U_t^b)$.  This is a degenerate case of the full deep BSDE
method — the $Z$-gradient that \citet{han2018} make the centrepiece
of their method plays no role — and a continuous-inventory
diffusion reformulation, in which $q_t$ follows a
stochastic-approximation SDE with $\sigma_q \neq 0$ and
$Z_t = \sigma_q\,\partial_q V$ is learned genuinely, is a natural
extension of the present work.  Taken together, the deep BSDE
literature up to 2023 provides every tool needed to solve the
Cont--Xiong problem — Han--E--Jentzen for the $Z$-gradient,
Han--Hu--Long for the MV distribution coupling, Wang et al.\ for the
jumps, Barles--Buckdahn--Pardoux for the analytic foundation — but
no existing study combines all three in the context of a realistic
market microstructure model.

But ``the tools exist'' is not the same as ``the method works
out of the box''.  Practical deep BSDE solvers for mean-field games
exhibit a cluster of failure modes that the optimistic original
literature does not document.  Three are particularly acute.
\emph{BatchNorm erasure}: when population-level statistics (such as
an aggregate competitor-quote summary $m_t$) are broadcast across a
minibatch, batch normalisation collapses their variance to zero and
the network learns a distribution-independent policy — the very
feature that mean-field theory was supposed to capture is silently
discarded.  \emph{Representation collapse under DeepSets}: the
permutation-invariant pooling operators routinely used to process
opponent states can, under symmetric initialisations, produce
degenerate representations that ignore competitor-specific
information.  \emph{Generator bypass}: when the network has enough
capacity, it can satisfy the terminal-loss
$|Y_T - g(q_T)|^2 \to 0$ with a $U$-profile that is internally
inconsistent with the first-order-condition-derived optimal quotes,
so that the solver converges to a self-consistent but non-Nash
fixed point — a pathology that a naive MV-fictitious-play BSDEJ
solver can exhibit.  \citet{germain2022} already catalogue
vanishing gradients, terminal-scale sensitivity, and $Y_0$
basin-locking; the three failure modes listed above are specific
to the \emph{mean-field} extension and to the
\emph{jump-augmented} solver, and are not systematically
addressed in the existing literature.  These failures share a
common structure: each of them eliminates or corrupts the
policy's dependence on the population distribution --- precisely
the feature that distinguishes a mean-field problem from a
single-agent one --- so that the network satisfies the training
loss without actually representing the mean-field interaction.
Documenting and repairing them is an open problem that this
thesis takes up.

% ============================================================
\section{Empirical Results and Financial Applications}\label{sec:lit-empirical}

The translation of deep BSDE and MV-FBSDE methods into finance has
proceeded along several parallel streams, each reporting its own
difficulties and partial successes.
\citet{germain2022} provides a critical survey of NN-based
algorithms for stochastic control and PDEs in finance, documenting
systematic training instabilities: vanishing gradients deep in the
time discretisation, sensitivity to the terminal-loss scale, and a
tendency for the $Y_0$ estimator to lock onto incorrect basins of
attraction when the terminal condition is not smooth.  Their paper
is a useful antidote to the optimism of the original deep BSDE
literature and motivates the careful regularisation and
warm-start procedures that any practical deep-BSDE solver for a
financial MFG must employ.  Rigorous convergence rates
for deep-NN stochastic control on finite horizons were obtained by
\citet{hure2020}, who combined deep RL with dynamic
programming and proved error bounds depending on the regularity of
the value function.

The pivot point of the empirical literature, for the purposes of the
present thesis, is \citet{han2022}.  By giving a deep-learning
algorithm for McKean--Vlasov FBSDEs whose convergence is free of the
curse of dimensionality under monotonicity, Han, Hu, and Long
unlocked the mean-field direction for deep BSDE.  Their algorithm
alternates between a fictitious-play outer loop, which updates the
empirical measure of the population, and an inner deep-BSDE solve,
which treats the measure as exogenous.  Convergence is established
in terms of the number of fictitious-play iterations and the
supervised-learning error in the inner BSDE.  The numerical
experiments focus on an ergodic linear-quadratic MFG and a systemic
risk model; the authors do not treat jump-driven games or discrete
execution, leaving a clear gap that the Cont--Xiong application
fills.

The Cont--Xiong model itself \citep{cont2024} is the thesis's
financial application blueprint.  Cont and Xiong model $N$ competing
market makers in a continuous-time dealer market as a stochastic
differential game of intensity control.  Each market maker posts bid
and ask quotes, and the execution probability in response to a
random order flow depends on the maker's own quote relative to a
softmax aggregate of competitors' quotes.  Nash equilibrium is
characterised by a coupled system of HJB equations (their
equation~28), whose existence is established by Brouwer's fixed-point
theorem, and whose computation Cont and Xiong reduce to a
fictitious-play iteration that alternates policy evaluation (a
linear Bellman solve) with a best-response step.  The Pareto-optimum
— which represents joint-objective collusion — is computed by a
parallel iteration on the joint objective.  Critically, Cont and
Xiong then show via decentralised multi-agent deep deterministic
policy gradient (MADDPG) that competing market-making algorithms,
learning independently from their own trade outcomes without any
shared information, converge to spreads strictly above the Nash
equilibrium, a phenomenon they term \emph{tacit collusion}.  This
empirical finding has direct regulatory implications and
reframes the MADDPG experiment as a detection mechanism: the gap
between learned and Nash spreads is a measurable signature of
algorithmic collusion even in the absence of explicit coordination.

A natural question follows.  If Cont and Xiong have already solved
their own model empirically via MADDPG, what does a BSDE-based
treatment add?  Three things.  \emph{Equilibrium grounding}: the
BSDEJ solution directly targets the Nash fixed point of the coupled
HJB system, not a behavioural attractor of a learning dynamic; the
recovered object is an equilibrium by construction.
\emph{Interpretability}: the learned $Z$ and $U$ processes have
direct economic meaning --- the sensitivity of the value function
to inventory and the jump in value at each execution event,
respectively --- whereas MADDPG's actor and critic networks are
opaque.  \emph{Convergence structure}: Han--Hu--Long-style
fictitious-play convergence, subject to the monotonicity caveat
noted in \S\ref{sec:lit-num}, comes with quantitative error
bounds; MADDPG offers none.  The two methods are therefore
complementary: MADDPG reveals what algorithms will actually do in
practice, while BSDEJ establishes the benchmark from which they
deviate.  The \emph{gap} between the two --- quantified as
$|S_{\mathrm{MADDPG}} - S_{\mathrm{Nash}}|$ --- is the natural
empirical object of interest for any BSDE-based treatment of the
game.

Related MARL work on realistic LOBs has produced complementary
results.  \citet{karpe2020} train multi-agent RL agents for
optimal execution in the ABIDES simulator, showing that model-free
RL can recover quasi-Almgren--Chriss policies when microstructure
signals are available; their setup is close to the Cont--Xiong
environment but lacks the stochastic-game structure.
\citet{shi2024} take a different angle, proposing a hybrid
neural stochastic agent-based model in which a pre-trained neural
point-process captures the background order-flow behaviour of
uncontrolled traders, and an agent-based layer on top introduces
strategic decisions.  Both of these papers highlight a common
empirical issue — generalisation across market regimes is fragile —
that motivates the use of principled BSDE-based solvers, whose
training only requires knowledge of the model coefficients rather
than historical data.

What the empirical literature lacks, as of late 2025, is a deep
BSDE-with-jumps solution of the Cont--Xiong MFG and a systematic
finite-$N$ to mean-field comparison.  The closest existing attempt
is the Han--Hu--Long algorithm applied to non-jump MFGs, combined
with the Cont--Xiong Bellman-plus-MADDPG study on the finance side.
Bridging these two is the technical aim of this thesis.

% ============================================================
\section{Conclusion and Open Problems}\label{sec:lit-conclusion}

Three very recent papers define the research frontier and help to
scope the contribution of the present thesis.
\citet{hu2025} extend the MV-FBSDE machinery to
environments with common noise by parameterising the measure flow
through path signatures rather than empirical measures.  Their deep
signature approach combines signature features — a tractable
compression of the history of a stochastic process — with neural
networks that operate on the signature rather than on the full path,
and establishes convergence in terms of fictitious-play iterations
under supervised-learning error control.  This is the natural
mean-field generalisation of common-noise MFGs
\citep{carmona2016, carmona2018} to the
numerical regime and offers a promising route for future extensions
of the present thesis to market making under systemic shocks.

\citet{navarro2025} propose a unified mean-field
framework that bridges informal bargaining markets and formal LOBs by
modelling both as interacting-particle systems on a one-dimensional
price lattice.  Their paper supports the central narrative of this
thesis — that dealer markets are natural mean-field environments —
and suggests directions for mapping the Cont--Xiong solver onto
fully general LOB topologies.  Finally, \citet{cheridito2025}
present ABIDES-MARL, a GPU-accelerated multi-agent reinforcement
learning environment built on top of ABIDES for studying endogenous
price formation in LOBs.  ABIDES-MARL addresses the scalability
limitations of earlier LOB simulators by exploiting parallelism for
both agent policy updates and environment simulation, providing a
platform on which the tacit-collusion experiments of
\citet{cont2024} can in principle be extended to hundreds
of agents.

Taken together, the literature gives a consistent diagnosis and a
narrow research gap.  BSDE theory and MFG theory are both mature,
and their probabilistic synthesis via MV-FBSDEs is well understood
\citep{pardoux1990, carmona2018,
buckdahn2009viscosity, chassagneux2014}.  Deep-learning
numerics for the Brownian-only case are well-studied, with both
empirical validation \citep{ehanjentzen2017, han2018,
beck2019} and partial convergence theory
\citep{carmona2019, hure2020}, and have recently
been extended to the curse-of-dimensionality-free MV setting
\citep{han2022}.  Financial applications to multi-agent
market making are only beginning: \citet{cont2024} provides
the stochastic-game model and a decisive empirical finding on tacit
collusion, and emerging frontiers include common-noise signatures
\citep{hu2025}, mean-field LOB embeddings
\citep{navarro2025}, and GPU-scale MARL environments
\citep{cheridito2025}.  What no existing paper provides is, to
state the central open question of this review in a single
sentence: a validated BSDE-with-jumps solution of the
Cont--Xiong game, an explanation of why naive deep BSDE methods
fail in that setting, or a quantitative operationalisation of
the MADDPG--Nash spread gap as a detector statistic for
algorithmic collusion.  These three absences motivate a single
coherent research arc --- \emph{making the Cont--Xiong
equilibrium computable, ensuring that the computation is
correct, and extracting economically meaningful structure from
the resulting solution} --- which this thesis takes up through
the following questions.  (i) Can a deep BSDEJ solver recover
the Cont--Xiong Nash equilibrium, and can the $O(1/\sqrt{N})$
propagation-of-chaos rate \citep{buckdahn2009limit} be verified
empirically on this specific model?  (ii) Which failure modes
(BatchNorm erasure, representation collapse, generator bypass)
derail naive applications of deep BSDE to mean-field games, and
how can they be systematically diagnosed and repaired?  (iii)
Can the spread gap $|S_{\mathrm{MADDPG}} - S_{\mathrm{Nash}}|$
be elevated from an empirical observation in \citet{cont2024}
to a measurement tool for regulators of algorithmic markets?
Taken together, these three questions define the scope of the
thesis that this review sets up.
